\documentclass{article}
\usepackage[T1]{fontenc}
\usepackage{lmodern}
\usepackage{graphicx}
\usepackage{amsmath,amssymb}
\usepackage[a4paper,margin=2cm]{geometry}
\usepackage[absolute,overlay]{textpos}
\setlength{\TPHorizModule}{1mm}
\setlength{\TPVertModule}{1mm}
\usepackage[indonesian]{babel}
\usepackage{hyperref}
\setlength{\emergencystretch}{2em}
% unit: Halaman 38

\begin{document}

% === Halaman 38 (python/native) ===

Penggunaan Kurva Eliptik di dalam Kriptografi

• Bagian inti dari sistem kriptografi kunci-publik yang melibatkan kurva eliptik 

adalah grup eliptik (himpunan titik-titik pada kurva eliptik dan sebuah operasi 

biner +).

• Operasi matematika yang mendasari:

     -  Jika RSA mempunyai operasi perpangkatan sebagai operasi matematika yang 

mendasarinya, maka

     -  ECC memiliki operasi perkalian titik (kP) 

38

*) Sumber bahan: Debdeep Mukhopadhyay, Elliptic Curve Cryptography ,

 Dept of Computer Sc and Engg IIT Madras

\end{document}
